The starting point is disclosure credibility. Accounting and finance research does not treat disclosure as a frictionless release of facts; it treats disclosure as communication shaped by information asymmetry, proprietary costs, reputational concerns, and strategic incentives (Healy \& Palepu, 2001; Verrecchia, 2001; Beyer et al., 2010). The key implication is more specific. When the activity being described is costly to implement and difficult to verify, the credibility of the language depends on whether lower-quality senders can cheaply mimic higher-quality ones (Spence, 1973; Seidmann \& Winter, 1997). AI disclosure fits that environment because filing language can expand more quickly than observable implementation.

The washing literature gives the same credibility problem empirical content. In environmental and sustainability settings, public claims can diverge from realized behavior when firms face strong incentives to appear aligned with valued goals (Delmas \& Burbano, 2011; Kim \& Lyon, 2015; Lyon \& Montgomery, 2015). Marquis et al.~(2016) show how selective disclosure can create the appearance of transparency despite weaker underlying records, while Hummel and Schlick (2016) emphasize that disclosure quality matters more than disclosure volume. Baker et al.~(2024) extend this logic to diversity disclosure, where claims can outpace later workforce outcomes. They motivate a test of whether valued disclosure aligns with implementation and later realization, beyond whether such language appears.

Recent AI-disclosure papers bring the credibility question closer to the present setting. Basnet et al.~(2025) provide the closest published starting point on the language side by distinguishing actionable, speculative, and irrelevant AI disclosure in 10-K business descriptions and showing that investors respond differently across those categories. Bandyopadhyay et al.~(2023) are closest on the market side: they show that broader AI narratives are associated with subsequent stock mispricing, especially among firms outside the largest size group. Related working papers study AI disclosure, disclosure--implementation gaps, and overstatement more directly (Li, 2025; Donelson et al., 2025; Barrios et al., 2025). They leave open how disclosure composition, contemporaneous implementation, and later realization can be linked without treating them as the same object. PatentMismatch addresses that gap by using filing-level disclosure composition to identify low-credibility AI narratives and relate them to realization, scrutiny, incentives, and market outcomes.

Textual-analysis research reaches a similar lesson from the measurement side: corporate language is informative only when it is measured in a way that fits the institutional setting. Forward-looking filing language predicts later fundamentals (Li, 2010), earnings-press-release language adds incremental information (Davis et al., 2012), and results are sensitive to the underlying vocabulary design (Loughran \& McDonald, 2011, 2016). Tone and wording can also be managed strategically rather than passively reflecting fundamentals (Huang et al., 2014). Narrative disclosure is not unreliable simply because it can be managed. Its empirical content depends on whether the measure isolates the relevant linguistic dimension rather than collapsing disclosure into a broad mention count. The distinction matters in an AI setting, where generic references, aspirational language, and implemented uses can all appear in the same filing while conveying very different information.

The innovation literature points in the same direction. Bellstam et al.~(2021) show that richer textual measures capture meaningful differences in firm-level innovation and forecast subsequent patenting and performance. Kim and Valentine (2023) show that public disclosure improves how markets interpret innovation-related assets. These studies point to a two-sided validation approach. Filing language can contain information about later technological activity, but measurement choices determine how much of that information is recovered. The narrower question is whether disclosure composition isolates filings less aligned with contemporaneous implementation and subsequent technological realization.
